\section{An analytic criterion for the Riemann hypothesis}\label{sec:criterion}

This section reduces the Riemann hypothesis to an upper bound for a smoothed prime sum on arbitrarily large power windows. The bound is one-sided, and the windows may be separated by arbitrarily large gaps. A zero to the right of the proposed exponent gives a pole of the Mellin transform. That pole forces positive Taylor coefficients at bounded gaps, whereas the window bound makes every coefficient in a long interval too small.

\subsection{The smoothed prime sum}\label{subsec:smoothed-prime-sum}

Let $\Lambda$ denote the von Mangoldt function and let $\gamma$ be Euler's constant. For $x\geq 1$, put
\begin{equation}\label{eq:prime-sum}
G(x)
=
\sum_{n\leq x}\Lambda(n)\left(\frac{x}{n}-1\right)
-x\log x+(1+\gamma)x.
\end{equation}
The last two terms remove the singularity at $s=1$ from the Mellin transform. The poles that remain in the right half-plane will detect the zeros of $\zeta$.

We first record the growth and regularity of $G$. Mertens' elementary estimate~\cite[Theorem~0.20]{vaughan-mertens} gives
\[
\sum_{n\leq x}\frac{\Lambda(n)}{n}
=
\log x+O(1)
\qquad(x\geq 2).
\]
The implied constant is absolute.

The function $G$ is continuous: the summand of index $n$ vanishes when $x=n$. Away from the integers,
\[
G'(x)
=
\sum_{n\leq x}\frac{\Lambda(n)}{n}-\log x+\gamma,
\]
so $G'$ is bounded by an absolute constant. Partitioning a finite interval at its integer points and applying the mean value theorem on each piece gives an absolute constant $K$ such that
\[
|G(x)-G(y)|\leq K|x-y|
\qquad(x,y\geq 1).
\]
Since $G(1)=1+\gamma$, there is also an absolute constant $B>0$ with
\begin{equation}\label{eq:linear-growth}
|G(x)|\leq Bx
\qquad(x\geq 1).
\end{equation}

For $\operatorname{Re}s>1$, the Euler product gives the absolutely convergent identity~\cite[Eq.~(27.4.12)]{dlmf-number-theory}
\[
-\frac{\zeta'(s)}{\zeta(s)}
=
\sum_{n\geq 1}\frac{\Lambda(n)}{n^s}.
\]
For every integer $n\geq 1$,
\[
\int_n^\infty
\left(\frac{x}{n}-1\right)x^{-s-1}\,dx
=
\frac{n^{-s}}{s(s-1)}.
\]
The interchange of summation and integration is justified absolutely. Indeed, if $\sigma=\operatorname{Re}s>1$, the sum of the corresponding absolute integrals is bounded by
\[
\frac{1}{\sigma(\sigma-1)}
\sum_{n\geq 1}\frac{\Lambda(n)}{n^\sigma}<\infty.
\]
Also,
\[
\int_1^\infty x\log x\,x^{-s-1}\,dx
=
\frac{1}{(s-1)^2},
\qquad
\int_1^\infty x\,x^{-s-1}\,dx
=
\frac{1}{s-1}.
\]
Consequently
\begin{equation}\label{eq:mellin-transform}
\begin{aligned}
F(s)
&:=\int_1^\infty G(x)x^{-s-1}\,dx\\
&=
-\frac{\zeta'(s)}{\zeta(s)}\frac{1}{s(s-1)}
-\frac{1}{(s-1)^2}
+\frac{1+\gamma}{s-1}
\qquad(\operatorname{Re}s>1).
\end{aligned}
\end{equation}

The meromorphic continuation and Laurent expansion of $\zeta$~\cite[Eq.~(25.2.4)]{dlmf-zeta} give
\[
-\frac{\zeta'(s)}{\zeta(s)}
=
\frac{1}{s-1}-\gamma+O(s-1).
\]
Since
\[
\frac{1}{s(s-1)}
=
\frac{1}{s-1}-1+O(s-1),
\]
the first term on the right of \eqref{eq:mellin-transform} has principal part
\[
\frac{1}{(s-1)^2}-\frac{1+\gamma}{s-1}.
\]
The two correction terms cancel this principal part. Thus $F$ extends meromorphically to $\mathbb C$ and is holomorphic at $s=1$.

Let $\rho$ be a zero of $\zeta$ with $\operatorname{Re}\rho>0$, of multiplicity $m_\rho$. Writing
\[
\zeta(s)=(s-\rho)^{m_\rho}h(s),
\qquad h(\rho)\ne 0,
\]
shows that $F$ has a simple pole at $\rho$, with residue
\begin{equation}\label{eq:pole-residue}
R_\rho
:=
\operatorname*{Res}_{s=\rho}F(s)
=
-\frac{m_\rho}{\rho(\rho-1)}.
\end{equation}
In particular, the residue is nonzero even when the zero is multiple. These are all the poles of $F$ in $\operatorname{Re}s>0$.

There are no positive real poles. For real $0<t<1$, the alternating representation of $\zeta$~\cite[Eq.~(25.2.3)]{dlmf-zeta} gives
\[
(1-2^{1-t})\zeta(t)
=
\sum_{j\geq 1}
\left((2j-1)^{-t}-(2j)^{-t}\right)>0.
\]
Since $1-2^{1-t}<0$, it follows that $\zeta(t)<0$. The point $s=1$ is removable for $F$, and $F$ is holomorphic in $\operatorname{Re}s>1$ by its integral representation. Finally, because $G$ is real,
\[
F(\overline{s})=\overline{F(s)}.
\]
Its poles therefore occur in conjugate pairs, with conjugate residues.

\subsection{Three lemmas}\label{subsec:lemmas}

The first lemma bounds the portions of an exponential integral outside a prescribed interval. Put
\[
I(t)=t-1-\log t
\qquad(t>0).
\]
Then $I(t)\geq 0$, with equality only at $t=1$. Moreover,
\[
I'(t)=1-\frac1t,
\]
so $I$ decreases on $(0,1)$ and increases on $(1,\infty)$.

\begin{lemma}\label{thm:exponential-tails}
Let $n\geq 1$ be an integer, let $a,U>0$, and put $z=aU/n$. If $z\leq 1$, then
\[
\int_0^U e^{-au}\frac{u^{n-1}}{(n-1)!}\,du
\leq
a^{-n}e^{-nI(z)}.
\]
If $z\geq 1$, then
\[
\int_U^\infty e^{-au}\frac{u^{n-1}}{(n-1)!}\,du
\leq
a^{-n}e^{-nI(z)}.
\]
\end{lemma}

\begin{proof}
Put $b=n/U$. When $z\leq 1$, we have $b\geq a$, and for $0\leq u\leq U$,
\[
e^{-au}
=
e^{-bu}e^{(b-a)u}
\leq
e^{(b-a)U}e^{-bu}.
\]
Hence
\[
\int_0^U e^{-au}\frac{u^{n-1}}{(n-1)!}\,du
\leq
e^{(b-a)U}
\int_0^\infty e^{-bu}\frac{u^{n-1}}{(n-1)!}\,du
=
e^{(b-a)U}b^{-n}.
\]
Since $aU=nz$ and $b=a/z$, the last expression equals
\[
a^{-n}\exp\!\bigl(n(1-z+\log z)\bigr)
=
a^{-n}e^{-nI(z)}.
\]

When $z\geq 1$, we have $b\leq a$, and the same comparison holds for $u\geq U$:
\[
e^{-au}
=
e^{-bu}e^{(b-a)u}
\leq
e^{(b-a)U}e^{-bu}.
\]
Integration over $[U,\infty)$ gives the second inequality.
\end{proof}

The next lemma ensures that the contribution of a conjugate pair of poles is positive at bounded gaps.

\begin{lemma}\label{thm:cosine-blocks}
Let $0<\varphi<\pi/2$, and let $\alpha\in\mathbb R$. Every block of
\[
N_\varphi=\left\lceil\frac{2\pi}{\varphi}\right\rceil
\]
consecutive integers contains an integer $j$ such that
\[
\cos(\alpha+j\varphi)\geq \frac12.
\]
\end{lemma}

\begin{proof}
The intervals on which the cosine is less than $1/2$ have length $4\pi/3$. They are separated by intervals of length $2\pi/3$ on which the cosine is at least $1/2$. Since the successive angles advance by
\[
\varphi<\frac{\pi}{2}<\frac{2\pi}{3},
\]
they cannot pass from one of the former intervals to the next without entering a latter interval.

If all $N_\varphi$ angles had cosine less than $1/2$, they would therefore lie in one interval of length $4\pi/3$. But their span is
\[
(N_\varphi-1)\varphi
\geq 2\pi-\varphi
>
\frac{4\pi}{3}.
\]
This is impossible.
\end{proof}

The last lemma selects a conjugate pair that determines the leading Taylor coefficients. It uses one specified pole to the right of the desired boundary; it does not require a pole of maximal real part.

\begin{lemma}\label{thm:nearest-poles}
Let
\[
P\subset\{s\in\mathbb C:\operatorname{Re}s\leq 1\}
\]
be a closed discrete set stable under complex conjugation, and suppose that $P$ contains no positive real point. Let $0\leq\theta<1$, and suppose that
\[
\rho_0=\beta_0+i\gamma_0\in P,
\qquad \beta_0>\theta.
\]
Put
\[
\theta'=\frac{\theta+\beta_0}{2}.
\]
For arbitrarily large real $\sigma>1$, the points of $P$ nearest to $\sigma$ form one conjugate pair $\rho,\overline{\rho}$, with $\operatorname{Im}\rho>0$. Their common distance $r$ satisfies
\[
\sigma-1<r<\sigma-\theta'.
\]
In particular, $r/(\sigma-1)$ tends to $1$ as $\sigma$ tends to infinity through these choices.
\end{lemma}

\begin{proof}
For all sufficiently large real $\sigma$,
\[
\begin{aligned}
(\sigma-\theta')^2-|\sigma-\rho_0|^2
&=
2\sigma(\beta_0-\theta')
+\theta'^2-|\rho_0|^2\\
&>0.
\end{aligned}
\]
The distance from $\sigma$ to $P$ is attained, because a bounded closed disc contains only finitely many points of $P$. Thus
\[
r\leq|\sigma-\rho_0|<\sigma-\theta'.
\]
Every nearest point $p$ consequently has $\operatorname{Re}p>\theta'$: otherwise
\[
|\sigma-p|
\geq \sigma-\operatorname{Re}p
\geq \sigma-\theta'.
\]
Such a point is nonreal, since $\theta'>0$.

For distinct $p,q\in P$, equality of their distances from a real $\sigma$ is equivalent to
\[
2\sigma(\operatorname{Re}q-\operatorname{Re}p)
=
|q|^2-|p|^2.
\]
If their real parts differ, this equation holds for at most one value of $\sigma$. If their real parts agree, equality of their absolute values forces $q=\overline p$. Since $P$ is countable, we may exclude a countable set of real centers and arrange that the nearest points form one conjugate pair.

Every point of $P$ has real part at most $1$, and $1\notin P$. Hence its distance from $\sigma>1$ is strictly greater than $\sigma-1$. We have proved
\[
1<
\frac{r}{\sigma-1}
<
\frac{\sigma-\theta'}{\sigma-1},
\]
whose right side tends to $1$.
\end{proof}

\subsection{The window criterion}\label{subsec:window-criterion}

\begin{theorem}\label{thm:window-criterion}
Let $0\leq\theta<1$. Suppose that there are constants $\lambda>1$ and $C>0$ such that, for arbitrarily large $Y$,
\begin{equation}\label{eq:window-bound}
G(x)\leq Cx^\theta
\qquad(Y\leq x\leq Y^\lambda).
\end{equation}
Then $\zeta$ has no zeros in the half-plane $\operatorname{Re}s>\theta$.
\end{theorem}

Here $C$ and $\lambda$ are fixed before the windows are chosen. The hypothesis means that for every $Y_0>1$, there is a $Y\geq Y_0$ for which \eqref{eq:window-bound} holds throughout the indicated interval.

The proof compares two expressions for the Taylor coefficients of $F$ at a fixed real point $\sigma>1$. The nearest conjugate pair of poles makes some of these coefficients positive and large. To bound the same coefficients from above, we split their integral representations at $\log Y$ and $\lambda\log Y$. After applying \eqref{eq:linear-growth}, the outer integrals have kernel
\[
e^{-(\sigma-1)u}\frac{u^{n-1}}{(n-1)!}.
\]
Lemma~\ref{thm:exponential-tails} makes both outer integrals exponentially small when $n/(\sigma-1)$ lies strictly inside the logarithmic window. The upper bound for $G$ controls the middle integral.

\begin{proof}
Suppose that $\zeta$ has a zero $\rho_0$ with
\[
\beta_0:=\operatorname{Re}\rho_0>\theta.
\]
By \eqref{eq:pole-residue}, this is a pole of $F$. Let $P$ be the set of all poles of $F$, and put
\[
\theta'=\frac{\theta+\beta_0}{2}.
\]

\emph{The base point.} Choose real numbers
\[
1<c_1<c_2<\lambda,
\]
and put
\[
d=
\min\left\{
I(1/c_1),\,
I(\lambda/c_2)
\right\}>0.
\]
The positivity follows from $1/c_1<1<\lambda/c_2$.

The pole set $P$ satisfies the hypotheses of Lemma~\ref{thm:nearest-poles}. We may therefore choose a real $\sigma>1$, with nearest poles $\rho,\overline{\rho}$, such that, on writing
\[
a=\sigma-1,
\qquad
r=|\sigma-\rho|,
\]
we have
\[
a<r<\sigma-\theta',
\qquad
\log\frac ra<\frac d2.
\]
Fix this choice, and put
\[
q=\frac{r}{\sigma-\theta}.
\]
Then $0<q<1$. All these quantities remain fixed as $Y$ varies.

\emph{The coefficients from the poles.} For an integer $n\geq 1$, define
\begin{equation}\label{eq:coefficient-integral}
T_{n-1}
=
\frac{(-1)^{n-1}}{(n-1)!}F^{(n-1)}(\sigma)
=
\int_0^\infty
G(e^u)e^{-\sigma u}
\frac{u^{n-1}}{(n-1)!}\,du.
\end{equation}
Differentiation under the integral is justified by \eqref{eq:linear-growth}.

Let $R_\rho\ne 0$ be the residue of $F$ at $\rho$. Subtracting the two principal parts gives
\[
Q(s)
=
F(s)
-\frac{R_\rho}{s-\rho}
-\frac{\overline{R_\rho}}{s-\overline{\rho}}.
\]
Because $\rho,\overline{\rho}$ are the only nearest poles and $P$ is discrete, there is a radius $R_*>r$ such that $Q$ is holomorphic on a neighborhood of the closed disc
\[
|s-\sigma|\leq R_*.
\]
Cauchy's estimate gives
\[
\frac{|Q^{(n-1)}(\sigma)|}{(n-1)!}
\leq M R_*^{1-n},
\qquad
M=\max_{|s-\sigma|=R_*}|Q(s)|.
\]

Write
\[
\sigma-\rho=re^{-i\varphi},
\qquad
R_\rho=|R_\rho|e^{i\alpha},
\]
where $0<\varphi<\pi/2$. Differentiating the principal parts in \eqref{eq:coefficient-integral}, we obtain
\begin{equation}\label{eq:coefficient-poles}
T_{n-1}
=
2|R_\rho|r^{-n}\cos(\alpha+n\varphi)+E_n,
\qquad
|E_n|\leq M R_*^{1-n}.
\end{equation}
Since $r<R_*$,
\[
r^n|E_n|
\leq
M R_*\left(\frac{r}{R_*}\right)^n,
\qquad
\lim_{n\to\infty}r^n|E_n|=0.
\]

Call an integer $n$ good if
\[
\cos(\alpha+n\varphi)\geq \frac12.
\]
Lemma~\ref{thm:cosine-blocks} gives a fixed integer $N_\varphi$ such that every block of $N_\varphi$ consecutive integers contains a good integer. It follows from \eqref{eq:coefficient-poles} that there is an $n_0$ such that every good $n\geq n_0$ satisfies
\begin{equation}\label{eq:coefficient-lower-bound}
T_{n-1}
\geq
\frac{|R_\rho|}{2}\,r^{-n}.
\end{equation}

\emph{The coefficients from a window.} Let $Y>1$ be one of the values in the hypothesis, and put $L=\log Y$. Consider integers $n$ in the interval
\begin{equation}\label{eq:index-window}
c_1aL\leq n\leq c_2aL.
\end{equation}
For these integers,
\[
\frac{aL}{n}\leq\frac1{c_1}<1,
\qquad
\frac{a\lambda L}{n}\geq\frac{\lambda}{c_2}>1.
\]
The monotonicity of $I$ gives
\[
I(aL/n)\geq I(1/c_1)\geq d,
\qquad
I(a\lambda L/n)\geq I(\lambda/c_2)\geq d.
\]
Applying Lemma~\ref{thm:exponential-tails} at the two ends of the window,
\begin{equation}\label{eq:tail-estimate}
\begin{aligned}
&\int_0^L e^{-au}\frac{u^{n-1}}{(n-1)!}\,du
+
\int_{\lambda L}^\infty
e^{-au}\frac{u^{n-1}}{(n-1)!}\,du\\
&\hspace{4em}\leq
2a^{-n}e^{-dn}
=
2r^{-n}\exp\!\left(n\log\frac ra-dn\right)
\leq
2r^{-n}e^{-dn/2}.
\end{aligned}
\end{equation}
On these outer intervals, \eqref{eq:linear-growth} gives
\[
|G(e^u)|e^{-\sigma u}\leq Be^{-au}.
\]

On the middle interval, \eqref{eq:window-bound} gives
\[
G(e^u)e^{-\sigma u}
\leq
Ce^{-(\sigma-\theta)u}.
\]
Since $C>0$,
\[
\begin{aligned}
\int_L^{\lambda L}
G(e^u)e^{-\sigma u}
\frac{u^{n-1}}{(n-1)!}\,du
&\leq
C\int_0^\infty
e^{-(\sigma-\theta)u}
\frac{u^{n-1}}{(n-1)!}\,du\\
&=
C(\sigma-\theta)^{-n}
=
Cq^n r^{-n}.
\end{aligned}
\]
Together with \eqref{eq:tail-estimate}, this proves the upper bound
\begin{equation}\label{eq:coefficient-upper-bound}
T_{n-1}
\leq
r^{-n}\left(2Be^{-dn/2}+Cq^n\right)
\qquad(c_1aL\leq n\leq c_2aL).
\end{equation}

\emph{The contradiction.} The length of the interval in \eqref{eq:index-window} satisfies
\[
(c_2-c_1)a\log Y \xrightarrow{Y\to\infty} \infty.
\]
For every sufficiently large admissible $Y$, it therefore contains $N_\varphi$ consecutive integers, all at least $n_0$. Choose a good integer $n$ among them. Comparing \eqref{eq:coefficient-lower-bound} and \eqref{eq:coefficient-upper-bound} gives
\[
\frac{|R_\rho|}{2}
\leq
2Be^{-dn/2}+Cq^n.
\]
The bound $n\geq c_1a\log Y$ makes $n$ arbitrarily large. Since $d>0$ and $0<q<1$, the right side tends to zero as $n\to\infty$, contradicting $R_\rho\ne 0$.
\end{proof}

In contrapositive form, Theorem~\ref{thm:window-criterion} says that a zero with real part greater than $\theta$ forces the following behavior: for every fixed $\lambda>1$ and $C>0$, every sufficiently large window $[Y,Y^\lambda]$ contains a point $x$ with
\[
G(x)>Cx^\theta.
\]
This is the obstruction that the arithmetic estimate must exclude.

\subsection{The estimate sufficient for the Riemann hypothesis}\label{subsec:rh-estimate}

The arithmetic argument will allow a small loss in the exponent. The relevant requirement is that the coefficient of that loss remain fixed as the exponent approaches $1/2$.

\begin{corollary}\label{thm:rh-criterion}
Let $c>0$ be fixed. Suppose that, for every sufficiently small $\varepsilon>0$, there are constants
\[
\lambda_\varepsilon>1,
\qquad
C_\varepsilon>0
\]
such that, for arbitrarily large $Y$,
\begin{equation}\label{eq:rh-window-bound}
G(x)
\leq
C_\varepsilon x^{1/2+c\varepsilon}
\qquad
(Y\leq x\leq Y^{\lambda_\varepsilon}).
\end{equation}
Then the Riemann hypothesis holds.
\end{corollary}

\begin{proof}
Suppose that $\zeta$ has a zero $\rho$ with
\[
\operatorname{Re}\rho>\frac12.
\]
Choose $\varepsilon>0$ small enough that the hypothesis applies and
\[
\frac12+c\varepsilon
<
\min\{\operatorname{Re}\rho,1\}.
\]
Theorem~\ref{thm:window-criterion}, with
\[
\theta=\frac12+c\varepsilon,
\qquad
\lambda=\lambda_\varepsilon,
\qquad
C=C_\varepsilon,
\]
excludes this zero. Thus $\zeta$ has no zeros with real part greater than $1/2$.

The functional equation~\cite[Eqs.~(25.4.3)--(25.4.4)]{dlmf-zeta}
\[
\xi(s)=\xi(1-s),
\qquad
\xi(s)=\frac12s(s-1)\pi^{-s/2}\Gamma(s/2)\zeta(s),
\]
reflects every nontrivial zero to a nontrivial zero on the other side of the line $\operatorname{Re}s=1/2$. Hence no nontrivial zero has real part less than $1/2$, and all nontrivial zeros lie on that line.
\end{proof}

For each fixed $\varepsilon$, the constants in \eqref{eq:rh-window-bound} must work throughout an unbounded family of windows. The constants, the starting threshold, and the windows themselves may depend on $\varepsilon$. The coefficient $c$ is fixed before $\varepsilon$ is chosen.

It therefore suffices to construct, for every sufficiently small $\varepsilon>0$, arbitrarily large windows $[Y,Y^{3/2}]$ on which
\[
G(x)\ll_\varepsilon x^{1/2+c\varepsilon}
\]
with one fixed $c$. The remaining sections of Part I reduce this estimate to a bound for the excess attached to an arithmetic chain.
